% Compilar no MacTeX:
%   latexmk -xelatex main.tex
% Limpar ficheiros auxiliares:
%   latexmk -c
\documentclass{istdissertacao}

\addbibresource{referencias.bib}

% -----------------------------------------------------------------------------
% DADOS DA DISSERTAÇÃO
% -----------------------------------------------------------------------------
\titulodissertação{Título da Dissertação}
\subtitulodissertação{Subtítulo}
\autordissertação{Nome completo do Candidato}
\graudissertação{Mestre}
\cursodissertação{Nome do Curso}
\orientadoresdissertação{Professor Doutor Nome Completo}
\presidentejuri{Professor Doutor Nome Completo}
\vogaisjuri{Professor Doutor Nome Completo\\Doutor Nome Completo}
\datadissertação{Setembro de 2026}
% \imagemcapa{figuras/imagem-capa.pdf} % altura máxima: 5 cm
% \logoacademia{figuras/logo-academia.pdf}
% \versaodefinitiva % descomentar para remover “Documento Provisório”

\begin{document}
\selectlanguage{portuguese}

\capadissertação
\declaracao

\frontmatterIST

\chapter*{Agradecimentos}
\addcontentsline{toc}{chapter}{Agradecimentos}
% Facultativo.
Os agradecimentos devem indicar de forma específica o contributo de cada pessoa ou organização.

\chapter*{Resumo}
\addcontentsline{toc}{chapter}{Resumo}
% Máximo: 250 palavras; sem referências bibliográficas.
Apresente aqui uma descrição sumária do trabalho realizado, da metodologia, dos principais resultados e das conclusões.

\vspace{1em}
\textbf{Palavras-chave:} palavra-chave 1; palavra-chave 2; palavra-chave 3; palavra-chave 4.

\cleardoublepage
\selectlanguage{english}
\chapter*{Abstract}
\addcontentsline{toc}{chapter}{Abstract}
% Maximum: 250 words; do not cite references.
Write here a concise summary of the work, methodology, main results and conclusions.

\vspace{1em}
\textbf{Keywords:} keyword 1; keyword 2; keyword 3; keyword 4.
\selectlanguage{portuguese}

\tableofcontents
\cleardoublepage
\listoftables
\addcontentsline{toc}{chapter}{Lista de Quadros}
\cleardoublepage
\listoffigures
\addcontentsline{toc}{chapter}{Lista de Figuras}

\chapter*{Lista de abreviaturas e símbolos}
\addcontentsline{toc}{chapter}{Lista de abreviaturas e símbolos}
\begin{tabularx}{\textwidth}{@{}lX@{}}
$A$ & Área da secção transversal\\
$E$ & Módulo de elasticidade\\
$L$ & Comprimento\\
$\alpha$ & Coeficiente a definir\\
\end{tabularx}

\mainmatterIST
\chapter{Introduction}
\label{cap:introducao}

\section{Context and motivation}
\label{sec:motivacao}

The Large Hadron Collider (LHC) at CERN delivers the highest energy proton-proton
collisions ever produced in a laboratory and, with them, an enormous sample of
strongly interacting final states. The majority of these collisions are governed
by Quantum Chromodynamics (QCD), the gauge theory of the strong interaction.
The study and comprehension of these collisions are not only a prerequisite for
suppressing backgrounds in searches for new phenomena but also the structure of
the hadronic final state can be examined and compared to theory directly, at scales
where the coupling is weak enough for perturbation theory to be applied and yet the
dynamics still remain rich.

The most prominent objects in these final states are jets. Jets are collimated sprays
of hadrons that emerge from the fragmentation of high energy quarks and gluons.
A jet is not a fundamental object but instead a definition. A clustering algorithm
specifies which detector inputs are grouped together and a recombination scheme
specifies how momentum is assigned to a result. This defines a jet, from which an
interesting and important question then arises. What does the radiation inside the jets look like?
This domain is called jet substructure, which over the past decade has grown from
a collection of tagging tools into a precision probe of QCD dynamics in its own right. 

In order to organise and look at the internal structure of a jet, one can use the Lund
jet plane \cite{Dreyer:2018nbf}. The Lund jet plane is a two dimensional map of a jet's
branching history on which soft, collinear and non perturbative regions of QCD radiation
occupy distinct territories. This map records where the jet's internal emissions occur, but
it sets aside an intriguing piece of information, the orientation of each splitting
relative to the one before. This orientation is not random. Through the linear polarisation
of intermediate gluons,
successive splittings are correlated in their relative azimuthal angle, producing
a characteristic modulation in the distribution of that angle. This is a genuine quantum
interference effect which survives only because the intermediate gluon is described by
an amplitude rather than a probability, and it is absent from any picture that treats successive
emissions as independent classical branchings.

These collinear spin correlations are well understood theoretically and are implemented in
modern parton shower algorithms, yet establishing them cleanly in collision data has
proved rather difficult. The reason is that the effect carries opposite signs in different
splitting channels. The $g \to q\bar{q}$ and $g \to gg$ channels contribute with opposite
sign, so a measurement that sums all channels sees the two contributions partially cancel.
This interference is not absent from such an inclusive measurement. It is diluted
and the apparent size reflects the accidental flavour composition of the sample rather than the
strength of the underlying physics.

This is the obstacle the present work addresses. A meaningful measurement of the
effect requires knowing which splitting produced each emision. However, the flavour
of an individual splitting inside a jet is not directly observable and must be inferred.
Working entirely on public CERN Open Data \cite{CERNOpenData}, more precisely with
CMS Open Data \cite{CMSOpenDataGuide}, this dissertation builds a reproducible pipeline
from the released samples through Lund declustering to the azimuthal spin observable,
and introduces a flavour selection that isolates the $g \to q\bar{q}$ channel. The
resulting modulation is reported not as a single number but as a function of channel
purity achieved by that selection, a
presentation that separates the physics of the interference from the performance
of the selection and allows the magnitude of the effect to be read off in the
limit of a pure sample.

The Lund declustering angle $\dpsi$ is not the only way to reach this effect.
The same $\cos 2\psi$ modulation also imprints itself on the flow of energy
within a jet, and can be accessed through a three-point energy correlator, which
probes the relative azimuthal orientation of energy deposited by triplets of
particles without reference to any clustering history. The two observables are
complementary: one reads the effect from the reconstructed branching sequence,
the other from the energy flow directly, and each carries a different exposure to
the choices and imperfections of the measurement. This work develops both, 
treating the Lund observable as the primary measurement and the energy correlator
as an independent cross-check on the same underlying physics.

The fact that this analysis rests on public data is not incidental.
The CERN Open Data Portal~\cite{CERNOpenData} has made it possible for work
outside the experimental collaborations to produce meaningful results and to
carry out genuinely new measurements. The present work is intended as a
contribution in that spirit. It builds an end to end pipeline from a single
CMS dataset to the final measurements, establishing along the way a workflow
for accessing, handling and measuring with CMS Open Data.

\section{Objectives}
\label{sec:objetivos}

The main aim of this dissertation is to measure collinear spin correlations in QCD
jets using public LHC data and to establish the flavour selection that such a
measurement requires in order to be meaningful. For the purpose of achieving that goal,
the work is organised around a sequence of objectives. The first is to build a reproducible 
pipeline over CERN Open Data samples capable
of extracting the necessary variables to begin a jet substructure analysis. This means
that the data samples must contain jet constituent information, thus the data format
chosen must retain that level of detail, ruling out the more compact samples in favour
of one that preserves per-constituent quantities.

With that pipeline in place, the second objective is to construct the primary
and secondary Lund jet planes and to verify that they reproduce the known
features of QCD radiation, since this agreement is what establishes that the
declustering is sound before any finer observable is built on top of it. 
Following that, the third objective is to implement the azimuthal observable $\dpsi$
between the hardest primary and secondary branchings, and to study its
distribution both with and without grooming, recovering the inclusive result
before attempting to sharpen it.

The fourth objective is the central one, and the point at which this work departs
from a purely inclusive analysis. Because the spin-interference signal carries
opposite signs in the $g \to q\bar{q}$ and $g \to gg$ channels, an inclusive
measurement dilutes it. Isolating a single channel is therefore not a refinement
but a precondition for measuring the effect at its true strength. This objective
is thus to develop a flavour selection capable of enriching the $g \to q\bar{q}$
channel, and to quantify both how well it performs and how pure a sample it can
reach in principle. The fifth objective builds directly on it. It is to measure
$\langle \cos 2\psi_{12} \rangle$ as a function of the channel purity delivered
by that selection and to establish the value characteristic
of a pure $g \to q\bar{q}$ sample. 

Finally, the work aims to assess the systematic
uncertainties affecting this measurement and to identify the dominant limitations
of the approach, so that the extensions it opens up can be set out on solid ground.

\section{Thesis outline}
\label{sec:estrutura}

Chapter \ref{cap:estado-da-arte} develops the theoretical material on which the
measurement rests, such as the elements of QCD relevant to jet formation, the
definition and reconstruction of jets and the algorithms used to build and
decluster them, the Lund jet plane, and the origin of collinear spin
correlations together with the two observables through which this work accesses
them. Chapter \ref{cap:metodos} turns to the data and the analysis, and covers the CERN
Open Data resources and the CMS detector, the data format and
dataset used, and the pipeline that runs from the released samples through Lund
declustering to the azimuthal observable, the flavour selection, and the
treatment of systematic uncertainties. Chapter \ref{cap:resultados} presents the
validation of the pipeline and the measurement itself, and discusses the results
in the light of the expectations set out in
Chapter~\ref{cap:estado-da-arte}. Finally,
Chapter~\ref{cap:conclusoes} returns to the objectives of
Section~\ref{sec:objetivos}, states what has been established and with what
confidence, and outlines the extensions that the present framework makes
possible.


\chapter{Theoretical Background}
\label{cap:estado-da-arte}

\section{Quantum Chromodynamics}
\label{sec:qcd}

Quantim Chromodynamics (QCD) is the quantum field theory of the strong interaction,
a non-Abelian gauge theory built on the colour symmetry group $SU(3)$ \cite{Ellis:1996mzs}.
Its fundamental degrees of freedom are quarks, which carry colour charge, and gluons,
the gauge bosons that mediate the interaction. The non-Abelian character of the theory
has a decisive consequence absent from Quantum Electrodynamics. The gluons themselves
carry colour, and therefore couple directly to one another. To a large extent, the
dynamics of the strong interaction are the dynamics of this gluon self-coupling, and
the two properties of QCD that govern everything in this chapter both follow from it.

\subsection{Asymptotic freedom and confinement}

The first property is asymptotic freedom. The strong coupling $\alpha_s$ decreases
as the momentum transfer of an interaction increases, so that at the high scales 
probed in LHC collisions the quarks and gluons interact weakly and perturbation
theory applies \cite{Gross:1973id,Politzer:1973fx}. The running of the coupling
is captured at leading order by

\begin{equation}
    \alpha_s(Q^2) = \frac{\alpha_s(\mu^2)}
    {1 + \alpha_s(\mu^2)\,\dfrac{33 - 2 n_f}{12\pi}\,\ln\!\dfrac{Q^2}{\mu^2}},
    \label{eq:running}
\end{equation}

where $Q$ is the scale of the process, $\mu$ a reference scale, and $n_f$ the
number of active quark flavours. The positive sign of the coefficient
$33 - 2n_f$, which holds for any physical number of flavours, is what makes the
coupling fall at large $Q$.  It is a direct imprint of the gluon self-interaction,
and its discovery established that a peturbative treatment of hard processes is
justified. It is convenient to trade the reference scale for a single dimensionful
constant $\Lambda_{\mathrm{QCD}}$, the scale at which the peturbative coupling formally
diverges, of an order of a few hundered MeV. In terms of it the coupling is small for
$Q \gg \Lambda_{\mathrm{QCD}}$ and grows strong as $Q$ approaches it.

This growth is the second property called confinement. At low momentum transfer
$\alpha_s$ becomes large and colour-charged particles cannot be isolated. A
quark or gluon produced in a hard scattering does not propagate freely to a detector,
because as it separates from other colour charges, the energy stored in the colour field
between them grows with distance until it becomes energetically favourable to create
a new quark-antiquark pairs from the vacuum. The original parton is thereby
dressed into a spray of colour-neutral hadrons. This is why the strongly interacting
final state reaches the detector not as partons but as the collimated sprays introduced
in Chapter \ref{cap:introducao}, and why any comparison between a peturbative
calculation and a measurement must bridge the divide between the patronic and hadronic
descriptions.

\subsection{Soft and collinear radiation}

The bridge between the hard parton and the observed hadrons is built from
successive emissions, and the pattern of those emissions is far from uniform.
When a parton radiates a gluon, the emission probability at first order in the
coupling takes the schematic form

\begin{equation}
  d\mathcal{P} \;\propto\; \frac{2\alpha_s C}{\pi}\,
  \frac{d\theta}{\theta}\,\frac{dE}{E},
  \label{eq:soft-collinear}
\end{equation}

where $E$ and $\theta$ are the energy and emission angle of the radiated gluon
and $C$ is a colour factor. The two factors carry the characteristic structure of
QCD radiation. The factor $d\theta/\theta$ diverges as $\theta \to 0$ meaning that
emmissions are enhanced when the gluon is radiated close to the direction of its emmiter.
This is known as the collinear divergence. The factir $dE/E$ diverges as $E \to 0$
meaning that emissions are enchanced when the gluon is soft. This is known
as the soft divergence. Together they mean that a hard parton radiates predominantly
softly and at small angles, and that number of such emissions is,
in a precise logarithmic sense, large.

The collinear structure is described more completely by the 
Dokshitzer–Gribov–Lipatov–Altarelli–Parisi (DGLAP) splitting
functions \cite{Ellis:1996mzs}, which give the probability that a parton of one type splits into two,
as a function of the momentum fraction $z$ carried by one of the products. For
the three splittings relevant to this work they read, at leading order,

\begin{align}
  P_{q \to qg}(z) &= C_F\,\frac{1 + z^2}{1 - z}, \label{eq:Pqq}\\[4pt]
  P_{g \to gg}(z) &= C_A\,
    \left[ \frac{z}{1-z} + \frac{1-z}{z} + z(1-z) \right], \label{eq:Pgg}\\[4pt]
  P_{g \to q\bar{q}}(z) &= T_R\,\left[ z^2 + (1-z)^2 \right], \label{eq:Pgq}
\end{align}

with colour factors $C_F = 4/3$, $C_A = 3$ and $T_R = 1/2$. The distinction
between the gluon splitting channels of Eqs. \eqref{eq:Pgg} and \eqref{eq:Pgq},
one producing a pair of gluons and the other a quark-antiquark pair, will
prove central. The two channels differ not only in these rates but in the spin
structure of their emissions, and it is that difference which the spin
observables of Section \ref{sec:spin} are designed to probe.

These divergences are not a pathology. They are integrable, and in any physical,
infrared-safe observable the divergent real emissions cancel against
corresponding virtual corrections, leaving a finite result. This cancellation
is the requirement that shapes the very definition of a jet in Section \ref{sec:jatos}.
Their logarithmic structure, however, is physical and measurable, and it is exactly
the structure that the Lund jet plane of Section \ref{sec:lund} is built to expose.
By mapping each emmision onto logarithmic axes of angle and transverse momentum,
it turns the double logarithmic enhancement of Eq. \eqref{eq:soft-collinear}
into an approximately uniform density, against which finer correlations stand out.

\subsection{Parton showers and hadronization}

Because a single hard parton typically radiates many times before hadronizing,
the emissions can be organised into a sequence orderred from the hardest,
widest-angle splittings down to the soft and collinear ones. This sequence
is what is known as a parton shower. Modern Monte Carlo event generators implement
this sequence probabilistically through the splitting functions of Eqs. 
\eqref{eq:Pqq} and \eqref{eq:Pgq}, evolving each parton down from the hard 
scale of the collision to a low scale of order $1\GeV$, at which the perturbative
description breaks down and confinement takes over. Below that scale the partons
are converted into hadrons by phenomenological hadronization models, tuned to data
rather than derived from first principles.

This picture, a hard, perturbatively calculable core, dressed by a shower of
increasingly soft and collinear emissions and finally hadronized, underlies
the whole of the analysis that follows. The observables studied in this work are
built to read the perturbative shower through the hadrons that reach the
detector, and for this reason the measurement is deliberately confined to a hard
kinematic regime, where the perturbative component dominates and non-perturbative
corrections are comparatively suppressed.

\section{Jets}
\label{sec:jatos}

As established in Section \ref{sec:qcd}, a coloured parton produced in a hard
scattering does not reach the detector directly. It radiates and hadronizes,
arriving as a collimated spray of hadrons. A jet is the operational device that
inverts this process approximately, grouping those hadrons back together so that
their combined momentum estimates that of the parton that produced them. A jet is
therefore not a physical object with a sharp definition but a construction, and
its properties depend on the prescription used to build it. This section sets out
that prescription: the requirement any jet definition must satisfy to be usable
in perturbative QCD, and the three algorithms, built on that requirement, on
which this work relies.

\subsection{Jet algorithms and infrared--collinear safety}
\label{sec:ircsafety}

Two ingredients specify how a jet is built. A jet algorithm decides which
inputs, from detector-level energy deposits, reconstructed particles, or, in a
calculation, partons, are grouped together into a jet. A recombination
scheme specifies how the momenta of those grouped inputs are combined into the
momentum of the jet. Throughout this work the inputs are the particle-flow
candidates described in Chapter \ref{cap:metodos}, and momenta are combined in the
$E$-scheme, in which four-momenta are simply added. This scheme is Lorentz invariant
and yields jets with non-zero mass.

Unfortunately, not every conceivable grouping prescription is admissible.
The requirement that disqualifies most of them is infrared and collinear (IRC) safety.
An observable is infrared safe if its value is unchanged by the addition of an
arbitrarily soft emission, and collinear safe if it is unchanged when one
of its inputs is replaced by two collinear inputs sharing its momentum.
A jet algorithm is IRC safe when the set of jets it returns is stable under both
these operations. A soft particle must not create or destroy a jet and a
collinear splitting must not change how particles are assigned to jets.

The reason why this requirement is not optional follows directly from
Section \ref{sec:qcd}. The emission probability of Eq. \eqref{eq:soft-collinear}
diverges in both the soft and the collinear limits. In any IRC-safe observable
these divergent real emissions cancel against equally divergent virtual
corrections, leaving a finite result. In an IRC-unsafe observable, that cancellation
is spoiled and a fixed-order peturbative prediction returns infinity rather than a number.
An IRC-unsafe jet definition is therefore not merely inconvenient but it also cannot
be compared with peturbative QCD at all. Since the purpose of measuring jets is to
confront the theory with data, IRC safety is a precondition and not a refinement.

This criterion has been decisive in the historical development of jet algorithms.
The earliest definitions were cone algorithms, which grouped particles falling 
within a fixed angular radius of a candidate axis. In their simplest, seed-based
forms these algorithms were not IRC safe because a soft particle could seed a new
candidate cone, or a collinear splitting could change which particles acted as seeds,
and either could alter the resulting jets. Much effort was devoted to repairing
them, and eventually producing seedless variants that restored collinear safety.
The alternative family, the sequential recombination algorithms, are IRC safe by
construction, and it is to this family that algorithm used in this work belongs.

It is worth making concrete how such a failure arises, since the same mechanism
will reappear when flavour is assigned to jets in Section \ref{sec:ifn}. Consider
a seed-based cone algorithm applied to an event with two hard particles whose
separation is slightly greater than the cone radius, so that each seeds its own
jet. The addition of a single soft particle between them may seed a third cone
that overlaps both, merging the two hard particles into one jet where before
there were two which, again, it is an infrared-unsafe change,
since an arbitrarily soft emission has altered the hard jets.
A collinear splitting produces the collinear-unsafe failure.
Replacing one hard seed by two collinear particles, each
carrying half its transverse momentum, may leave neither above the seed threshold,
so that a jet which existed before the splitting is absent after it. In both cases,
the divergent emissions of Eq. \eqref{eq:soft-collinear} feed directly into the
jets themselves, and the perturbative prediction diverges with them.

The sequential recombination algorithms introduced below evade both failures by
never relying on seeds or on fixed geometric cones. They build jets by repeatedly
merging inputs according to a distance measure, so that soft or collinear addition
changes only the order in which already nearby inputs are combined, never the final
assigment of hard momentum to jets. This structural difference is the reason 
the entire famiky is IRC safe, and it is what recommends it over the
cone algorithms for any measurement to be compared with theory.

A sequential recombination algorithm is defined by a distance measure between
each pair of inputs and between each input and the beam. The generalised-$\kt$
family \cite{Salam_2010} writes these as

\begin{align}
  d_{ij} &= \min\!\left( \pt{}_{,i}^{2p},\, \pt{}_{,j}^{2p} \right)
           \frac{\dR_{ij}^2}{R^2}, \label{eq:dij}\\[4pt]
  d_{iB} &= \pt{}_{,i}^{2p}, \label{eq:diB}
\end{align}

where $\dR_{ij}^2 = (y_i - y_j)^2 + (\phi_i - \phi_j)^2$ is the squared angular
separation of the inputs in rapidity $y$ and azimuth $\phi$, $R$ is the jet radius
parameter, and the integer $p$ selects the algorithm. At each step the algorithm
finds the smallest of all the distances. If it is a $d_{ij}$, the inputs $i$ and
$j$ are recombined into a single input. If it is a $d_{iB}$, the input $i$ is
declared a final jet and removed. The procedure repeats until no inputs remain.
The set of jets returned is IRC safe for any integer value of $p$, but the value
of $p$ changes profoundly how the clustering proceeds. The three choices $p=-1$,
$p=0$ and $p=+1$ give, respectively, the \antikt{}, Cambridge/Aachen and $\kt$
algorithms, of which the first two are the tools this work relies on and are
described in the following subsections.

Common to all three is the radius parameter $R$, which sets the angular reach of
the clustering through the ratio $\dR_{ij}^2/R^2$ in Eq. \eqref{eq:dij}. A smaller
$R$ yields narrower jets that resolve nearby structure but capture less of a single
parton's radiation. A larger $R$ collects more of that radiation at the price of
admitting more contamination from unrelated soft activity. The appropriate choice
depends on the use to which the jets are put, and the two algorithms this work
relies on are used with different values of $R$ for correspondingly different
purposes. This point is, again, taken up in Sections \ref{sec:antikt}
and \ref{sec:ca}.

\subsection{The anti-$\kt$ algorithm}
\label{sec:antikt}

The choice $p = -1$ in Eqs. \eqref{eq:dij} and \eqref{eq:diB} defines the \antikt{}
algorithm \cite{Cacciari_2008}. With this value the inter-particle distance becomes

\begin{equation}
  d_{ij} = \min\!\left( \frac{1}{\pt{}_{,i}^{2}},\, \frac{1}{\pt{}_{,j}^{2}} \right)
           \frac{\dR_{ij}^2}{R^2},
  \label{eq:antikt-distance}
\end{equation}

so that the distance scales with the inverse squared transverse momenta of
the inputs. The smallest distances, and therefore the first recombinations,
are those involving the hardest particles. The clustering is consequently driven
by hard inputs. A hard particle accretes the softer inputs around it before soft
inputs cluster among themselves.

The name of the algorithm reflects this inversion. In the $\kt$ algorithm, the
choice $p = 1$, the distance scales with the positive powers of transverse
momentum, so that soft inputs are clustered first and the jet is assembled from
the bottom up, following the inverse of the physical showering order. The
\antikt{} algorithm reverses this. Taking $p = -1$ turns the soft-first ordering
of $\kt$ into a hard-first one. It is, in effect, the $\kt$ algorithm run with the
role of soft and hard interchanged, and the prefix ``anti'' records precisely that
inversion. The two sit at opposite ends of the generalised family, with
Cambridge/Aachen, insensitive to transverse momentum altogether, between them.

This ordering has a striking geometric consequence. A hard particle with no other
hard particle within an angular distance $2R$ collects all softer inputs within a
distance $R$ of it, forming a jet that is exactly circular in the $(y,\phi)$
plane. Where two hard particles lie closer than $2R$, the boundary between their
jets is set by their relative transverse momenta, the harder one retaining the
larger share, but the jets remain regular in outline. Soft radiation, precisely
because it enters the clustering last, cannot deform these boundaries. It is
distributed among the hard jets without reshaping them.

A useful way to characterise this stability is through the notion of jet area, the
region of the $(y,\phi)$ plane from which a jet gathers its inputs. This area may
be probed operationally by adding a large number of infinitesimally soft
``ghost'' particles to an event and recording which jet each is assigned to. For
most algorithms the resulting areas are irregular and fluctuate from jet to jet,
because soft particles influence the clustering. However, for the \antikt{} algorithm,
where soft particles are clustered last and cannot reshape the hard jets, the
areas are close to the ideal value $\pi R^2$ and vary little. A stable, predictable
area is what makes the subtraction of pile-up and underlying-event contamination
tractable, since that contamination scales with the area over which each jet
collects it.

It is this regularity that has made the \antikt{} algorithm the standard jet
definition at the LHC. Because the jet boundary is stable against soft radiation,
the area over which a jet collects contamination from pile-up and the underlying
event is well defined and nearly constant, and the corresponding energy
corrections are straightforward to calibrate. The algorithm thus combines the IRC
safety guaranteed for the whole generalised-$\kt$ family with a jet shape that
is simple to treat experimentally, a combination that
its two siblings, whose boundaries are irregular, do not offer.

The prominence of the algorithm is also a matter of timing. When the LHC began
taking data, a jet definition was needed that was both IRC safe, so that
measurements could be compared with the increasingly precise perturbative
calculations then becoming available, and computationally fast enough to run on
the very large event samples the machine would produce. The \antikt{} algorithm,
together with its efficient implementation in the \textsc{FastJet}
package \cite{Cacciari_2012}, met both needs, and it was adopted as the baseline jet
definition by the LHC experiments, a role it has held since.

For these reasons the \antikt{} algorithm was used in this work to define
the jets, with radius parameter $R = 0.8$. Jets constructed in this way are
conventionally denoted as AK8. What the algorithm does not provide, however, is a
physically meaningful account of how the jet was assembled. Because its sequence is
ordered by hardness rather than by angle, the intermediate recombinations bear no
correspondence to the sequence of physical splittings in a parton shower, which is
ordered in angle. The clustering history of an \antikt{} jet is, in this sense, an
artefact of the algorithm rather than a record of the underlying radiation.
Recovering that history, which is what the substructure observables of this work
require, calls for a different member of the family, applied to the constituents
of the jet once it has been defined. That algorithm is the Cambridge/Aachen algorithm
and it is the subject of the next subsection.

\subsection{The Cambridge/Aachen algorithm}
\label{sec:ca}

\section{The Lund jet plane}
\label{sec:lund}

\section{Collinear spin correlations}
\label{sec:spin}
\chapter{Methods}
\label{cap:metodos}
\chapter{Resultados e Análise de Resultados}
Apresente os resultados de forma concisa e realize uma análise crítica, incluindo comparações com trabalhos anteriores, vantagens e limitações.

\chapter{Conclusões}
Sintetize o trabalho, relacione os resultados com os objetivos e apresente as principais recomendações e limitações.

\section{Desenvolvimentos futuros}
Indique possíveis extensões do trabalho.


\cleardoublepage
\printbibliography[heading=bibintoc,title={REFERÊNCIAS BIBLIOGRÁFICAS}]

\appendix
\chapter{Informação complementar}
Inclua nos anexos tabelas extensas, deduções, fichas técnicas, código ou outros elementos complementares.


\end{document}
